%%%%%%
%
% $Autor: Wings $
% $Datum: 2020-01-18 11:15:45Z $
% $Pfad: WuSt/Skript/Produktspezifikation/powerpoint/ImageProcessing.tex $
% $Version: 4620 $
%
%%%%%%


%todo Überschäumer entfernen


\chapter{Datenbasis}

Eine gute Datenbasis ist der Schlüssel für ein solides neuronales Netz. Wie aus Kapitel
4 bereits bekannt, erlernen neuronale Netze die Fähigkeit zu klassifizieren auf Grundlage
von Daten, die ihnen während des Trainings zur Verfügung gestellt werden. Je größer
und qualitativ hochwertiger die Datenbasis ist, desto stabiler und genauer kann das Netz
Klassifizierungen vornehmen. Dieses Kapitel befasst sich mit dem Thema, der Erstellung
einer guten Datenbasis.

\section{Aufbau eines Fotos}

Ein \ac{cnn} benötigt Eingangsdaten in Form einer Matrix, um trainiert werden zu können.
Ein Foto kann als solche Matrix beschrieben werden. Dabei wird das Bild in rasterförmig
angeordnete Bildpunkte, den sogenannten Pixeln, unterteilt. Jedem Pixel wird ein Farbwert 
zwischen 0 - 255 zugeordnet. Bei einem Schwarz-Weiß-Bild besitzt die Matrix die
Tiefe 1. Ein weißer Pixel hat den Wert 0 und ein schwarzer den Wert 255. Ein optisch
grauer Punkt hat einen Wert zwischen >0 und <255. Je heller der Pixel, desto kleiner
der Wert. Handelt es sich um ein Farbfoto so gibt es mehrere Schichten, die jeweils einen
Farbwert widerspiegeln. Die Bildtiefe ist dann >1. Die Größe der Matrix entspricht der
Größe des Bildes. \cite{Python:2020b}

\section{Vorbereitung der Daten}

Die Abbildung+\ref{DBSchema} zeigt einen Ablaufplan zum Erzeugen von Trainingsdaten. Nachfolgend
werden die einzelnen Punkte des Plans, anhand der ursprünglichen Datenbasis erläutert.
Erzeugt werden die Bilder an der Taktstraße 30 und den dort montierten Kameras der Firma 
Keyence. Diese erzeugen bei jeden Produktionstakt der Maschine ein Bild von einem
bestimmten Ausschnitt. Der Bildausschnitt ist so festgelegt, dass von jedem Werkzeug
der mittlere Teil fotografiert wird. Doppelformen bei denen links und rechts die gleichen
Zusatzfedern produziert werden, können mit einem Bild gleichzeitig geprüft werden. Die
Industriekameras die das Bild erzeugen sind an einen Sondercomputer der Firma Keyence
angeschlossen, welcher die Bilddateien über ein File Transfer Protocol (FTP)-Server an
einen Computer an der Taktstraße sendet. Dort wird das Bild in einem zuvor definierten
Ordner gespeichert.

Um den Computer vor äußeren Produktionsbedingungen und vor Fremdeinwirkungen zu
schützen, steht dieser in einem verschlossenen Metallschrank und ist mit einem Passwort
gesichert. Befugte Personen können mit entsprechendem Schlüssel und den Zugangsdaten
auf den Computer und die Bilder zugreifen. Da der Computer, aus sicherheitstechnischen
Gründen nicht mit dem Internet verbunden ist, wurde ein externes Speichermedium verwendet, 
um die Bilder aus dem Ordner zu entnehmen.

Bei einer Doppelform können sich insgesamt vier Zustände im Bild ergeben. Der erste
Zustand ist die vollständige Entnahme der Zusatzfedern aus beiden Formen, beziehungsweise
ein erfolgreicher Schließprozess und somit kein Überschäumer in einer Form. Dieser
Zustand wurde im nachfolgenden als \texttt{good good} definiert. Bei dem Zustand zwei und drei,
weist eine Seite eine unvollständige Entnahme beziehungsweise den Fehler
Überschäumer auf. Diese Kategorie wurden als good bad oder bad good bezeichnet, je nach dem auf
welcher Seite der Fehler auftritt. Der letzte Zustand ist eine fehlerhafte Entnahme auf
beiden Seiten oder der Fehler Überschäumer in beiden Formen. Dieser Zustand wurde als
\texttt{bad bad} festgelegt.

Für die Sortierung der Bilder auf dem Speichermedium wurde eine ruhige Arbeitsatmosphäre 
in einem Büro gewählt. In einem ausgewählten Verzeichnis auf dem Computer aus
dem Büro wurden vier Ordner der oben genannten Kategorien erstellt. Anschließend wurden 
die Fotos vom Speichermedium manuell gesichtet und in die entsprechenden Ordner
auf dem Computer einsortiert. Der symmetrische Aufbau des Bildes macht es möglich,
die Kategorien good good, good bad, bad good und bad bad auf die zwei Kategorien good
und bad zu reduzieren. Es wurden zwei neue Ordner mit dem Namen good und bad in
dem Verzeichnis angelegt. Mit Hilfe eines kurzen Programmabschnitt wurden die Bilder
aus den vier kategorisierten Ordnern geöffnet, in der Mitte geschnitten und in die Ordner
der Kategorien good oder bad als Kopie abgespeichert. Der Inhalt dieses Programms wird
im nachfolgenden Abschnitt erläutert.

Von jeder Kategorie wurde eine Testmenge von ca. 10\% zurückgehalten und in einen
separaten Ordner sortiert. Die Testmenge wurde manuelle ausgewählt, um einen repräsentativen 
Schnitt zu erhalten. Alternativ kann diese Aufgabe ebenfalls von einem
Programm übernommen werden. Als letzten Schritt wurden die verbleibenden Bilder aus
den Kategorien good und bad mit in einen großen Byte-Stream umgewandelt und gespeichert. 
Dieses Vorgehen wird als Serialisieren bezeichnet. \cite{Duden:2020}

\section{Programm Bilder zuschneiden und neu sortieren}


Das folgende Programm bearbeitet die Bilder aus den vier Kategorien good good, good bad,
bad good und bad bad, schneidet diese in der Mitte und speichert die Bildhälften in den
Ordner good oder bad. Das Programm wurde von Dr.-Ing. Ole Mathis Magens erstellt
und zur Verfügung gestellt.



Zu Beginn des Programms, werden die Bibliotheken aufgerufen, aus denen im nachfolgenden 
Programmcode Funktionen verwendetet werden. Um nicht die ganze Bibliothek in
das Programm laden zu müssen, können aus bestimmt Bibliotheken einzelne Funktionen
direkt aufgerufen werden. Aus der Bibliothek os wurde nur die Funktion \PYTHON{path} benötigt,
um auf Verzeichnisse zugreifen zu können.

Die Raw-Dateien liegen in einer sehr hohe Auflösung vor, das neuronale Netz kann zum Extrahieren 
der relevanten Informationen auch mit kleineren Auflösungen genaue Ergebnisse
erzielen. Daher wurde eine Ausgangsbreite (\PYTHON{output width}) und eine Ausgangshöhe 
(\PYTHON{output height}). als Variable definiert. In dem ersten Programm wurde die Ausgangsbreite 
und die Ausgangshöhe von einer Bildhälfte auf das Format 30x80 gesetzt.

\lstinputlisting%[firstline=300,lastline=500]
{../Python/JetsonNano/Zuschnitt.py}

Im nächsten Schritt bekommt die Variable files den Pfad zugeordnet in denen sich
die Bilder befinden. Mit dem Befehl glob.glob wird die Liste der Dateien mit ihrem
vollständigen Pfad zurückgegeben. Die Endung *.bmp gibt an in welchem Format die
Dateien sich befinden müssen, um berücksichtigt zu werden. [Fou20a] Es wird eine 
for-Schleife in Kombination mit einem range()-Befehl verwendet, um die nachfolgenden 
Befehle wiederholt über die gesamte Länge der im Ordner befindenden Bilder (len(files))
auszuführen. Mit dem open()-Befehl wird das Bild geöffnet und anschließend eine 
Transposition durchgeführt. Die von der Kamera gespeicherten Bilder sind in ihrer Ausrichtung
nach rechts gedreht. Der Befehl ROTATE 270 innerhalb der transpose()-Funktion dreht
das Bild um weitere 270 Grad nach rechts. Mit dem resize()-Befehl, wird die Größe des
Bildes geändert. Wie zu Beginn definiert, soll die Ausgangsgröße einer Bildhälfte 30x80
sein. Das ursprüngliche Bild besitzt die doppelte Bildbreite, da es sich um zwei Bildhälften
handelt. Der crop()-Befehl schneidet das Bild rechteckig zu. Die ersten beiden Werte innerhalb 
des Befehls definieren den Anfangspunkt (x 0 ,y 0 ) die folgenden zwei Werte definieren
den Endpunkt(x 1 ,y 1 ) des Zuschnitts. Das erste Bild geht von der unteren linken Ecke des
Bildes, bis zur oberen Hälfte. Die linke Bildhälfte wird anschließend im Ordner bad und
mit dem Zusatz left und den ursprünglichen Pfadnamen versehen. Danach wird die rechte 
Bildhälfte analog zugeschnitten und im Ordner good gespeichert. Mit dem Ausdruck
FLIP in dem transpose()-Befehl, wird das Bild nach rechts gespiegelt. Hierdurch wird eine
Vereinheitlichung der rechten und linken Seite erzeugt. Dieser Programmabschnitt wird
für die anderen Ordner good good, good bad und bad bad analog durchgeführt. Hierfür
muss das Verzeichnis entsprechend umgeändert werden. Derselbe Vorgang wurde für die
Bilder zur Erkennung von Überschäumern durchgeführt. In der Abbildung~\ref{} sind die
zugeschnittenen und gespiegelten Bilder dargestellt. Es ist zu erkennen, dass die Bilder
bis auf leicht veränderte Lichtverhältnisse identisch sind.

Nachdem dieses Programm durchlaufen wurde, befinden sich alle Ursprungsbilder als 
geschnittene Kopie in den Ordnern good und bad. In der Tabelle 6.1 wird die Anzahl der
Bilder in diesen Kategorien dargestellt, die für das erste Training verwendet wurde. Wie
sich anhand der Tabelle erkennen lässt, herrscht ein starkes Ungleichgewicht zwischen
den einzelnen Kategorien. Sowohl bei den Überschäumern und bei der Entnahme sind die
Bilder der Kategorie bad stark unterrepräsentiert.

\section{Bilden der Trainingsdaten}

Die für das Training benötigten Bilddateien befanden sich alle in kategorisierten Ordnern.
Um nicht jedes Bild einzeln in das neuronale Netz zu spielen, wurden alle Bilder für das
Training, in einem großen Byte-Stream serialisiert. Hierfür wird die pickle-Funktion aus
der gleichnamigen Bibliothek genutzt.

Der nachfolgende Programmcode zeigt die wichtigsten Inhalte zum Erzeugen der Trainingsdaten 
in reduzierter Form. Der vollständige Programmcode befindet sich im Anhang
und orientier sich an eine Vorlage aus dem Internet. \cite{sentdex:2018} Zu Beginn werden die 
notwendigen Bibliotheken aufgerufen. Wichtige Bibliotheken in diesem Programm sind das
cv2-Modul, die random-Bibliothek und die Bibliothek imblearn-over sampling welche ein
Datengleichgewicht zwischen unausgeglichenen Kategorien erzeugen kann.


\lstinputlisting%[firstline=300,lastline=500]
{../Python/JetsonNano/TrainingData.py}

In der Variablen \PYTHON{training\_data} werden die Inhalte der Bilder gespeichert. Mit der 
Funktion create trainings data wird dieser Inhalt erzeugt. Die erste for-Schleife dieser Funktion
erzeugt ein Zugriffspfad namens path, welcher einzelne Pfade zu den jeweiligen Kategorien 
herstellt. Mit der Variablen class num wird jeder Kategorie ein Index zugeordnet. Da
es zwei Kategorien gibt, bekommt die erste Kategorie good den Index 0 und die zweite
Kategorie bad den Index 1.

Die zweite for-Schleife liest die Bilder über die gesamte Länge der Ordnerinhalte und
speichert diese als Array in der Variable img arry. Der Befehl IMREAD Grayscale konvertiert
das Farbbild in ein Graustufenbild, wodurch die Bildtiefe von drei Ebenen auf eine reduziert 
wird. In der Variable new array wird der Inhalt des img array mit der geänderten
Bildtiefe gespeichert. Nun befinden sich die Bildinformationen aus beiden Kategorien in
der Variable new array. Zusammen mit den Informationen aus der Variable class num,
wird daraus der Trainingsdatensatz erstellt.

\begin{lstlisting}
random.shuffle( training_data )
\end{lstlisting}


Der Trainingsdatensatz war zu Beginn des Trainings noch sortiert. Das Bedeutet, dass die
ersten abgespeicherten Bilderinformationen aus der Kategorie good und die letzten aus
der Kategorie bad stammen. Um diese Sortierung zu durchbrechen wurde mit dem Befehl
random.shuffle gearbeitet. Diese Funktion schafft eine Durchmischung der Trainingsdaten.

\begin{lstlisting}
Xval = [ ]
yval = [ ]
X = [ ]
y = [ ]
ValNum =300
\end{lstlisting}

Die Variablen Xval, yval, X, Y und ValNum wurden neu definiert. Die Variable ValNum
gibt an, wie viele Bilder als Validierungsdaten zum Trainieren des Netzes verwendet wer-
den sollen. Der ursprüngliche Wert lag bei 300 Bilder.


\begin{lstlisting}
for features , label in training_data [0: ValNum - 1]:
    Xval.append ( features )
    yval.append ( label )
    
Xval = np.array (Xval).reshape( - 1 , IMG_HEIGHT , IMG_WIDTH , 1)

for features , label in training_data[ ValNum : - 1]:
    X.append( features )
    y.append( label )

X = np.array(X).reshape( - 1 , IMG_HEIGHT , IMG_WIDTH , 1)
\end{lstlisting}


Die for-Schleife hat die Funktion, die ersten 300 Positionen aus den Trainingsdaten zu
nehmen und die Bildinformationen (features) der Variable Xval hinzuzufügen. Die 
Information der zugehörigen Kategorie(lable) wird in der Variable yval gespeichert. Die
Variable Xval wurde anschließend als Numpy-Array der Größe 80x30 und der Tiefe 1
gespeichert. Analog hierzu wurden die verbleibenden Trainingsdaten als Variable X und
y gespeichert.

\begin{lstlisting}
sm = RandomOverSampler( random_state=0)
X , y = sm.fit_resample(X.reshape ( - 1 , IMG_HEIGHT*IMG_WIDTH ) , y)
X = X.reshape( - 1 , IMG_HEIGHT , IMG_WIDTH , 1)
\end{lstlisting}


Da in den einzelnen Kategorien ein starkes Ungleichgewicht der Bilderanzahl herrscht,
wurde ein Oversampling angewendet. Dazu wird die Funktion RandomOverSampling verwendet,
welche beliebige Bilder der unterrepräsentierten Kategorie vervielfältigt bis ein
Gleichgewicht zwischen den Kategorien herrscht.


\begin{lstlisting}
pickle_out = open("Xval.pickle" ,"wb")
pickle.dump(Xval , pickle_out )
pickle_out.close ()

pickle_out = open("yval.pickle" ,"wb")
pickle.dump(yval , pickle_out )
pickle_out.close ()

ickle_out = open("X.pickle" ,"wb")
pickle.dump(X , pickle_out )
pickle_out.close ()

pickle_out = open("y.pickle" ,"wb")
pickle.dump(y , pickle_out )
pickle_out.close ()
\end{lstlisting}

Die pickle-Funktion schreibt die zuvor erstellten Variablen Xval, yval, X und y in einen
Byte-Stream. Die Datei besitzt die Endung .pickle. Das wb in dem Programmabschnitt
steht für write bytes.

Nach diesem Vorgang war das Erstellen der Trainingsdaten abgeschlossen.